\documentclass[11pt,a4paper]{article}
\usepackage[utf8]{inputenc}
\usepackage[T1]{fontenc}
\usepackage[french]{babel}
\usepackage{amsmath,amssymb}
\usepackage{geometry}
\usepackage{enumitem}
\usepackage{booktabs}
\usepackage{array}
\usepackage{tabularx}
\usepackage{xcolor}
\usepackage{tcolorbox}
\usepackage{fancyhdr}
\usepackage{multicol}
\usepackage{tikz}
\usepackage{pgfplots}
\pgfplotsset{compat=1.18}
\usepackage{lastpage}

\tcbuselibrary{skins,breakable}

\geometry{margin=2cm, headheight=15pt}

% Couleurs
\definecolor{primary}{RGB}{31,78,121}
\definecolor{soft}{RGB}{230,239,247}
\definecolor{amberbg}{RGB}{255,244,224}
\definecolor{amberborder}{RGB}{240,192,112}
\definecolor{greenbg}{RGB}{233,246,234}
\definecolor{greenborder}{RGB}{182,224,184}
\definecolor{redbg}{RGB}{253,236,236}
\definecolor{redborder}{RGB}{242,179,171}
\definecolor{grayline}{RGB}{120,130,140}

% En-tête et pied de page
\pagestyle{fancy}
\fancyhf{}
\fancyhead[L]{\textcolor{primary}{\textbf{BTS BioALC — 2\textsuperscript{ème} année}}}
\fancyhead[R]{\textcolor{primary}{Fiche de travail — Dispersion}}
\fancyfoot[L]{\textcolor{grayline}{\small Nom : \dotfill}}
\fancyfoot[C]{\textcolor{grayline}{\small \thepage/\pageref{LastPage}}}
\fancyfoot[R]{\textcolor{grayline}{\small Semaine du 02/11/2026}}
\renewcommand{\headrulewidth}{0.4pt}
\renewcommand{\footrulewidth}{0.4pt}

% Boîtes
\newtcolorbox{consigne}{
  colback=soft, colframe=primary, boxrule=0.8pt, arc=4pt,
  title={\textbf{Mode d'emploi}}, fonttitle=\bfseries,
  coltitle=white, colbacktitle=primary,
  left=6pt, right=6pt, top=4pt, bottom=4pt,
  breakable
}

\newtcolorbox{contexte}{
  colback=amberbg, colframe=amberborder, boxrule=0.8pt, arc=4pt,
  title={\textbf{Contexte professionnel}}, fonttitle=\bfseries,
  coltitle=black, colbacktitle=amberborder,
  left=6pt, right=6pt, top=4pt, bottom=4pt,
  breakable
}

\newtcolorbox{rappel}[1][]{
  colback=greenbg, colframe=greenborder, boxrule=0.8pt, arc=4pt,
  fonttitle=\bfseries, coltitle=black, colbacktitle=greenborder,
  left=6pt, right=6pt, top=4pt, bottom=4pt,
  breakable, #1
}

\newtcolorbox{questionbox}[1][]{
  colback=white, colframe=primary, boxrule=0.8pt, arc=4pt,
  fonttitle=\bfseries, coltitle=white, colbacktitle=primary,
  left=6pt, right=6pt, top=4pt, bottom=4pt,
  breakable, #1
}

\newtcolorbox{autoeval}{
  colback=redbg, colframe=redborder, boxrule=0.8pt, arc=4pt,
  title={\textbf{Auto-évaluation}}, fonttitle=\bfseries,
  coltitle=black, colbacktitle=redborder,
  left=6pt, right=6pt, top=4pt, bottom=4pt,
  breakable
}

\newtcolorbox{evaluation}{
  colback=redbg, colframe=redborder, boxrule=0.8pt, arc=4pt,
  title={\textbf{Évaluation formative 1 — QCM}}, fonttitle=\bfseries,
  coltitle=black, colbacktitle=redborder,
  left=6pt, right=6pt, top=4pt, bottom=4pt,
  breakable
}

% Commandes utilitaires
\newcommand{\lignes}[1]{\vspace{0.15cm}\foreach \x in {1,...,#1}{\noindent\rule{\linewidth}{0.4pt}\par\vspace{0.35cm}}}
\newcommand{\case}{\raisebox{0.1ex}{\fbox{\rule{0pt}{1.1ex}\rule{1.1ex}{0pt}}}}
\newcommand{\carreau}{\raisebox{0.1ex}{\fbox{\rule{0pt}{1.4ex}\rule{1.4ex}{0pt}}}}

% Compteur pour les couples (i), (ii), (iii)
\newcommand{\couplei}{\textsf{(i)}}
\newcommand{\coupleii}{\textsf{(ii)}}
\newcommand{\coupleiii}{\textsf{(iii)}}

\title{\textcolor{primary}{\textbf{Fiche de travail — Caractéristiques de dispersion}}\\
\large Semaine du 2 au 6 novembre 2026 — BTS Bioanalyses en Laboratoire de Contrôle}
\author{}
\date{}

\begin{document}

\maketitle

\begin{consigne}
Complétez cette fiche directement pendant les deux séances. Les espaces libres sont prévus pour vos réponses, vos calculs et vos schémas. Cette fiche n'est pas notée : elle vous servira pour réviser et préparer l'évaluation formative de fin de semaine. En cas de difficulté, reportez-vous aux fiches méthodes rappelées à chaque partie.
\end{consigne}

\vspace{0.3cm}

\begin{contexte}
Toujours dans le cadre du contrôle qualité du laboratoire d'analyses médicales, le responsable vous demande maintenant d'évaluer la \textbf{dispersion} des mesures de taux de glucose obtenues sur 20 patients. Deux séries de mesures ont été réalisées :
\begin{itemize}[leftmargin=*]
    \item \textbf{Série A} — mesures de routine, un seul technicien ;
    \item \textbf{Série B} — mesures inter-laboratoires, plusieurs techniciens.
\end{itemize}
Objectif : déterminer si les deux séries sont \textbf{homogènes} et si la méthode de mesure est \textbf{fiable}.
\end{contexte}

\vspace{0.4cm}

% ============================================================
\section*{Partie 1 — Séance du lundi (50 min) : étendue, variance, écart-type}
% ============================================================

\subsection*{1.1. Rebrassage — Rappel de la semaine précédente}

\begin{questionbox}
\begin{enumerate}[leftmargin=*]
    \item Rappelez la définition de la \textbf{moyenne}, de la \textbf{médiane} et du \textbf{mode}.
    \lignes{2}
    \item Quelle est la différence entre un indicateur de \textbf{position} et un indicateur de \textbf{dispersion} ?
    \lignes{2}
\end{enumerate}
\end{questionbox}

\subsection*{1.2. Fiche méthode n\textsuperscript{o}3 — Étendue, variance, écart-type}

\begin{rappel}[title={\textbf{Fiche méthode — Indicateurs de dispersion}}]
\begin{tabularx}{\linewidth}{|p{2.6cm}|X|p{3.6cm}|}
\hline
\textbf{Indicateur} & \textbf{Définition / Formule} & \textbf{Interprétation} \\
\hline
\textbf{Étendue} & $E = x_{\max} - x_{\min}$ & Facile à calculer, mais très sensible aux valeurs extrêmes. \\
\hline
\textbf{Variance} & $V = \dfrac{1}{n}\displaystyle\sum_{i=1}^{n}(x_i - \bar{x})^2$ & Mesure l'écart moyen au carré par rapport à la moyenne. Unité : carré de celle des données. \\
\hline
\textbf{Écart-type} & $\sigma = \sqrt{V}$ & Même unité que les données. Plus $\sigma$ est grand, plus la série est dispersée. \\
\hline
\end{tabularx}

\vspace{0.2cm}
\textbf{Exemple type :} série $4, 5, 6, 8, 9, 10$.

Moyenne : $\bar{x} = \dfrac{4+5+6+8+9+10}{6} = \dfrac{42}{6} = 7$.

Variance : $V = \dfrac{(4-7)^2 + (5-7)^2 + (6-7)^2 + (8-7)^2 + (9-7)^2 + (10-7)^2}{6} = \dfrac{9+4+1+1+4+9}{6} = \dfrac{28}{6} \approx 4{,}67$.

Écart-type : $\sigma = \sqrt{4{,}67} \approx 2{,}16$.
\end{rappel}

\subsection*{1.3. Exercice 1 — Étendue et écart-type sur une série simple}

\begin{questionbox}
Soit la série de mesures de pH suivante obtenue sur 8 échantillons d'eau :

\begin{center}
\begin{tabular}{|c|c|c|c|c|c|c|c|}
\hline
7,1 & 7,4 & 6,9 & 7,3 & 7,0 & 7,5 & 6,8 & 7,2 \\
\hline
\end{tabular}
\end{center}

\textbf{Question 1.} Calculez l'\textbf{étendue} de cette série.

\lignes{2}

\textbf{Question 2.} Calculez la \textbf{moyenne} de la série.

\lignes{2}

\textbf{Question 3.} Complétez le tableau de calcul de la variance :

\vspace{0.2cm}
\begin{center}
\begin{tabular}{|c|c|c|}
\hline
\textbf{Valeur $x_i$} & \textbf{$x_i - \bar{x}$} & \textbf{$(x_i - \bar{x})^2$} \\
\hline
7,1 & & \\
\hline
7,4 & & \\
\hline
6,9 & & \\
\hline
7,3 & & \\
\hline
7,0 & & \\
\hline
7,5 & & \\
\hline
6,8 & & \\
\hline
7,2 & & \\
\hline
\textbf{Total} & & \\
\hline
\end{tabular}
\end{center}

\vspace{0.2cm}
Variance : $V = \dfrac{\dots\dots\dots}{\dots\dots\dots} = \dots\dots\dots$

\vspace{0.2cm}
Écart-type : $\sigma = \sqrt{V} = \dots\dots\dots$

\vspace{0.2cm}
\textbf{Question 4.} Le pH attendu est de 7,2 avec une tolérance d'écart-type $\sigma \leq 0{,}3$. La série est-elle conforme ?

\lignes{3}
\end{questionbox}

\subsection*{1.4. Exercice 2 — Comparaison de deux séries}

\begin{questionbox}
Deux séries de mesures de glucose (g/L) obtenues sur 10 patients :

\vspace{0.2cm}
\textbf{Série A :} 0,85 · 0,88 · 0,86 · 0,87 · 0,89 · 0,84 · 0,86 · 0,88 · 0,87 · 0,85

\textbf{Série B :} 0,75 · 0,82 · 0,95 · 0,79 · 0,88 · 0,92 · 0,81 · 0,86 · 0,78 · 0,94

\vspace{0.2cm}
\textbf{Question 1.} Calculez la moyenne et l'étendue de chaque série.

\lignes{3}

\textbf{Question 2.} Sans faire tous les calculs, laquelle des deux séries vous semble la plus dispersée ? Justifiez.

\lignes{3}

\textbf{Question 3.} Calculez l'écart-type de chaque série (vous pouvez utiliser le tableur).

\lignes{3}

\textbf{Question 4.} Quelle série est la plus fiable ? Que peut-on en déduire sur la méthode de mesure ?

\lignes{3}
\end{questionbox}

\subsection*{1.5. Synthèse de la séance 1}

\begin{questionbox}
\textbf{À retenir — complétez les phrases :}
\begin{itemize}[leftmargin=*]
    \item L'étendue est la différence entre \dotfill
    \item La variance est la moyenne des \dotfill
    \item L'écart-type est la \dotfill de la variance.
    \item Plus l'écart-type est grand, plus la série est \dotfill
\end{itemize}
\end{questionbox}

\newpage

% ============================================================
\section*{Partie 2 — Séance du jeudi (50 min) : quartiles et contrôle qualité}
% ============================================================

\subsection*{2.1. Rebrassage — Rappel de la séance 1}

\begin{questionbox}
\begin{enumerate}[leftmargin=*]
    \item Quelle est la différence entre l'étendue et l'écart-type ?
    \lignes{2}
    \item Pourquoi l'étendue est-elle sensible aux valeurs extrêmes ?
    \lignes{2}
\end{enumerate}
\end{questionbox}

\subsection*{2.2. Fiche méthode n\textsuperscript{o}4 — Quartiles et écart interquartile}

\begin{rappel}[title={\textbf{Fiche méthode — Quartiles}}]
\begin{tabularx}{\linewidth}{|p{2.6cm}|X|p{3.6cm}|}
\hline
\textbf{Indicateur} & \textbf{Définition / Méthode} & \textbf{Interprétation} \\
\hline
\textbf{$Q_1$} & Premier quartile : valeur telle que 25\,\% des données lui soient inférieures ou égales. & 1\textsuperscript{er} quart de la série. \\
\hline
\textbf{$Q_2$} & Deuxième quartile = \textbf{médiane}. & 50\,\% des données. \\
\hline
\textbf{$Q_3$} & Troisième quartile : valeur telle que 75\,\% des données lui soient inférieures ou égales. & 3\textsuperscript{e} quart de la série. \\
\hline
\textbf{IQR} & Écart interquartile : $IQR = Q_3 - Q_1$. & Mesure de dispersion \textbf{robuste} aux valeurs extrêmes. \\
\hline
\end{tabularx}

\vspace{0.2cm}
\textbf{Méthode de calcul :} trier les valeurs, couper la série en 4 groupes d'effectifs égaux, puis lire les valeurs aux frontières.

\textbf{Repérage des valeurs aberrantes :} une valeur est suspecte si elle sort de l'intervalle
\[
\bigl[\,Q_1 - 1{,}5 \times IQR\ ;\ Q_3 + 1{,}5 \times IQR\,\bigr].
\]
\end{rappel}

\subsection*{2.3. Exercice 1 — Calcul des quartiles et de l'IQR}

\begin{questionbox}
Reprenez la série de pH de la séance 1 (8 valeurs triées) :

\begin{center}
\begin{tabular}{|c|c|c|c|c|c|c|c|}
\hline
6,8 & 6,9 & 7,0 & 7,1 & 7,2 & 7,3 & 7,4 & 7,5 \\
\hline
\end{tabular}
\end{center}

\textbf{Question 1.} Déterminez $Q_1$, $Q_2$ (médiane) et $Q_3$.

\lignes{3}

\textbf{Question 2.} Calculez l'écart interquartile $IQR = Q_3 - Q_1$.

\lignes{2}

\textbf{Question 3.} Déterminez l'intervalle de repérage des valeurs aberrantes $[Q_1 - 1{,}5 \times IQR\ ;\ Q_3 + 1{,}5 \times IQR]$.

\lignes{3}

\textbf{Question 4.} Une des mesures est $6{,}0$. Est-elle aberrante ? Justifiez.

\lignes{3}
\end{questionbox}

\subsection*{2.4. Exercice 2 — Application au contrôle qualité}

\begin{questionbox}
Reprenez les deux séries A et B de la séance 1 (mesures de glucose).

\vspace{0.2cm}
\textbf{Question 1.} Complétez le tableau suivant (aide au tableur) :

\vspace{0.2cm}
\begin{center}
\begin{tabular}{|l|c|c|}
\hline
\textbf{Indicateur} & \textbf{Série A} & \textbf{Série B} \\
\hline
Moyenne $\bar{x}$ & & \\
\hline
Médiane $Q_2$ & & \\
\hline
$Q_1$ & & \\
\hline
$Q_3$ & & \\
\hline
$IQR$ & & \\
\hline
Écart-type $\sigma$ & & \\
\hline
\end{tabular}
\end{center}

\vspace{0.3cm}
\textbf{Question 2.} Pour chaque série, indiquez le couple d'indicateurs le plus pertinent parmi : \couplei{} (moyenne/variance), \coupleii{} (médiane/IQR), \coupleiii{} (médiane/mode). Justifiez en une phrase.

\lignes{3}

\textbf{Question 3.} Rédigez une conclusion de 4 à 5 lignes à destination du responsable qualité : quelle série est \textbf{homogène} ? La méthode de mesure est-elle \textbf{fiable} ? Y a-t-il des \textbf{valeurs aberrantes} à signaler ?

\lignes{5}
\end{questionbox}

\subsection*{2.5. Synthèse de la séance 2}

\begin{questionbox}
\textbf{À retenir — complétez :}
\begin{itemize}[leftmargin=*]
    \item $Q_1$ sépare les \dotfill\,\% des valeurs les plus faibles du reste.
    \item $Q_2$ est égal à la \dotfill
    \item $IQR = \dotfill$
    \item L'IQR est \textbf{robuste} car il n'est pas influencé par les \dotfill
    \item Une valeur est suspecte si elle sort de l'intervalle \dotfill
\end{itemize}
\end{questionbox}

\newpage

% ============================================================
\section*{Évaluation formative 1 — QCM statistique descriptive}
% ============================================================

\begin{evaluation}
\textbf{Durée : 10 minutes — sans note — correction immédiate.}

\vspace{0.3cm}
Pour chaque question, entourez la bonne réponse.

\vspace{0.4cm}
\textbf{Question 1.} La moyenne est un indicateur :

\hspace{1cm} A. de position \hspace{2.5cm} B. de dispersion \hspace{2.5cm} C. de forme

\vspace{0.4cm}
\textbf{Question 2.} L'écart-type mesure :

\hspace{1cm} A. la valeur centrale d'une série

\hspace{1cm} B. la dispersion des valeurs autour de la moyenne

\hspace{1cm} C. la valeur la plus fréquente

\vspace{0.4cm}
\textbf{Question 3.} L'unité de l'écart-type est :

\hspace{1cm} A. le carré de l'unité des données

\hspace{1cm} B. la même que celle des données

\hspace{1cm} C. sans unité

\vspace{0.4cm}
\textbf{Question 4.} Pour la série $4, 5, 5, 6, 7, 7, 7, 8, 9, 10$, la médiane vaut :

\hspace{1cm} A. 5 \hspace{2.5cm} B. 6{,}5 \hspace{2.5cm} C. 7

\vspace{0.4cm}
\textbf{Question 5.} Un écart important entre moyenne et médiane indique :

\hspace{1cm} A. une distribution symétrique

\hspace{1cm} B. la présence possible de valeurs extrêmes

\hspace{1cm} C. une erreur de calcul

\vspace{0.4cm}
\textbf{Question 6.} Le couple d'indicateurs adapté à une distribution asymétrique avec valeurs aberrantes est :

\hspace{1cm} A. (moyenne, variance)

\hspace{1cm} B. (médiane, interquartile)

\hspace{1cm} C. (médiane, mode)

\vspace{0.4cm}
\textbf{Question 7.} L'écart interquartile $IQR$ est égal à :

\hspace{1cm} A. $Q_2 - Q_1$ \hspace{2.5cm} B. $Q_3 - Q_1$ \hspace{2.5cm} C. $Q_3 - Q_2$

\vspace{0.4cm}
\textbf{Question 8.} Une valeur est considérée comme suspecte si elle sort de l'intervalle :

\hspace{1cm} A. $[\bar{x} - \sigma\ ;\ \bar{x} + \sigma]$

\hspace{1cm} B. $[Q_1 - 1{,}5 \times IQR\ ;\ Q_3 + 1{,}5 \times IQR]$

\hspace{1cm} C. $[x_{\min}\ ;\ x_{\max}]$

\vspace{0.4cm}
\textbf{Question 9.} Dans un contexte de contrôle qualité, la médiane est souvent préférée à la moyenne car :

\hspace{1cm} A. elle est plus facile à calculer

\hspace{1cm} B. elle est moins sensible aux valeurs aberrantes

\hspace{1cm} C. elle est toujours plus grande

\vspace{0.4cm}
\textbf{Question 10.} Une série a un écart-type très faible. Cela signifie que :

\hspace{1cm} A. les valeurs sont très dispersées autour de la moyenne

\hspace{1cm} B. les valeurs sont proches les unes des autres

\hspace{1cm} C. la moyenne est proche de zéro

\end{evaluation}

\newpage

% ============================================================
\section*{Bilan de la semaine}
% ============================================================

\begin{autoeval}
Cochez la case correspondant à votre niveau pour chaque compétence :

\vspace{0.3cm}
\begin{center}
\begin{tabularx}{\linewidth}{|X|c|c|c|}
\hline
\textbf{Compétence} & \textbf{Acquis} & \textbf{En cours} & \textbf{À revoir} \\
\hline
Je sais calculer une étendue. & \case & \case & \case \\
\hline
Je sais calculer une variance et un écart-type. & \case & \case & \case \\
\hline
Je sais déterminer $Q_1$, $Q_2$ et $Q_3$. & \case & \case & \case \\
\hline
Je sais calculer et interpréter l'écart interquartile. & \case & \case & \case \\
\hline
Je sais repérer une valeur aberrante. & \case & \case & \case \\
\hline
Je sais choisir le couple d'indicateurs adapté. & \case & \case & \case \\
\hline
Je sais rédiger une conclusion pour le contrôle qualité. & \case & \case & \case \\
\hline
\end{tabularx}
\end{center}

\vspace{0.3cm}
\textbf{Une question que je me pose encore :}

\lignes{2}
\end{autoeval}

\vspace{0.4cm}

\begin{contexte}
\textbf{Pour aller plus loin (facultatif, non noté) :} recherchez la notion de \textbf{carte de contrôle} (control chart) utilisée en contrôle qualité. Expliquez en 5 lignes comment la moyenne et l'écart-type permettent de détecter une dérive d'un procédé de mesure.

\lignes{5}
\end{contexte}

\vspace{0.5cm}

\begin{center}
\textit{Fiche de travail — BTS BioALC — 2\textsuperscript{ème} année — Semaine du 02/11/2026}\\
\textit{Document à conserver pour les révisions.}
\end{center}

\end{document}